% CarbonScope India — IEEE conference skeleton. Tables are generated: `make paper` (runs make_tables.py first).
\documentclass[conference]{IEEEtran}
\usepackage{booktabs,amsmath,graphicx,hyperref,siunitx}
\newcommand{\todo}[1]{\textcolor{red}{[TODO: #1]}}
\usepackage{xcolor}

\begin{document}
\title{CarbonScope India: Deep Panel Forecasting of State- and City-Level CO\textsubscript{2} Emissions for All 36 Indian States and Union Territories}
\author{\IEEEauthorblockN{\todo{Author names}}\IEEEauthorblockA{\todo{Affiliation}}}
\maketitle

\begin{abstract}
We forecast fossil CO\textsubscript{2} emissions 1, 3 and 5 years ahead for all 28 states and 8 union territories of India.
Because official state inventories are incomplete, targets are EDGAR 2025 gridded emissions area-weighted to post-2020
state boundaries; drivers combine CEA plant-level power data, MoSPI GSDP, Census population, NASA POWER climate and
OxCGRT lockdown stringency. With only $\sim$35 annual observations per state, we train a single global model with learned
state embeddings and compare 11 methods under a strict time split (train $\le$2014, validation 2015--2019, test 2020--2024).
State forecasts are made coherent with the national total via MinT reconciliation. \todo{one-sentence headline result from Table~\ref{tab:test}.}
\end{abstract}

\section{Introduction}
India's updated NDC and net-zero-by-2070 pledge are implemented largely through state policy (power procurement,
EV policies, industrial permitting), yet sub-national emission projections are scarce. Contributions:
\begin{enumerate}
\item A comparative deep-learning benchmark on all 36 Indian states/UTs under a leakage-safe protocol
      (to our knowledge the first at this coverage \todo{verify against literature}).
\item A state-embedding global panel design that makes deep models usable for small-N annual panels.
\item Hierarchical MinT reconciliation so state forecasts sum exactly to the national total; we report both
      unreconciled and reconciled accuracy.
\item A city-level proxy estimation framework (EDGAR grid buffers) with explicit estimation-method and confidence tiers.
\item (v2) The first India state$\times$sector deep-learning benchmark: 252 CO\textsubscript{2}e series (7 sectors,
      AR6 GWP-100) with a factorised state+sector embedding.
\item (v2) Two-level MinT reconciliation (national $\rightarrow$ \{states, sectors\} $\rightarrow$ cells) for emissions.
\item (v2) A cross-learning multi-task model that forecasts all sectors of a state in one pass.
\item (v2) A transparent scenario/counterfactual engine and an open, deterministic insight-generation framework (ATAR).
\end{enumerate}

\section{Data}
\textbf{Target.} EDGAR 2025 GHG CO\textsubscript{2} TOTALS, 0.1$^\circ$, 1990--2024, aggregated by exact cell/polygon
area overlay (EPSG:6933) to 36 state/UT polygons (Ladakh separate; DNH\&DD merged). These are \emph{EDGAR-downscaled
estimates}, not official inventories. \textbf{Validation source.} CEA CO\textsubscript{2} Baseline Database v22
(plant-level, FY2017--2025). \textbf{Drivers.} EDGAR sector grids (power, industry, road, cement, buildings);
MoSPI GSDP (2011-12 series) back-cast with World Bank real GDP; Census 2001/2011 population interpolated to 2023/24
estimates; NASA POWER monthly temperature/precipitation (CDD, JJAS rainfall); OxCGRT subnational stringency; policy
calendars (UJALA, PAT, JNNSM phases, state EV policies). All imputations are logged (\texttt{imputation\_report.csv}).

\section{Method}
Sliding windows of $W{=}10$ years predict log-growth $y_{t+h}=\log(1{+}c_{t+h})-\log(1{+}c_t)$ for $h{=}1..5$. Each
(sample, horizon) pair is assigned to a split by its \emph{target} year; scalers, winsorisation bounds and income tiers
are fitted on train years only (enforced by unit tests). Deep models (LSTM, GRU, TCN, Seq2Seq with Bahdanau attention,
TFT-lite with variable selection and multi-head attention, DeepAR-style Gaussian head) share one network across states
with an 8-d state embedding and static covariates (region, coastal, UT, income tier), trained with pinball loss weighted by
the inverse per-state MASE scale. MinT uses shrinkage covariance when residual rows exceed series, otherwise its WLS
(variance-scaled) form.

\section{Results}
\begin{table}[t]\centering\caption{Test (2020--2024) accuracy, all states. MASE scaled by train-period naive-1 MAE.
$^\dagger$ARIMA fitted on the 10 largest emitters only.}\label{tab:test}
\resizebox{\columnwidth}{!}{\begin{tabular}{lrrrrr}
\toprule
Model & MASE$_1$ & MASE$_3$ & MASE$_5$ & MinT MASE$_5$ & sMAPE$_5$ \\
\midrule
TFT-lite & 2.92 & 3.75 & \textbf{2.90} & 2.92 & 11.1 \\
Seq2Seq-Attn & 2.23 & 3.73 & 3.07 & 3.08 & 10.0 \\
LSTM + emb. & 1.72 & \textbf{3.25} & 3.45 & 3.48 & 11.5 \\
ARIMA$^\dagger$ & 1.75 & 3.58 & 3.91 & 3.91 & 11.5 \\
Naive (RW) & 2.18 & 3.84 & 4.02 & 3.66 & 12.8 \\
DeepAR-style & 1.98 & 4.41 & 4.51 & 4.46 & 14.7 \\
RW + drift & 1.82 & 4.08 & 5.13 & 5.12 & 15.7 \\
TCN + emb. & 1.93 & 4.26 & 6.28 & 6.25 & 18.7 \\
LightGBM (panel) & \textbf{1.72} & 3.90 & 6.64 & 6.63 & 19.9 \\
GRU + emb. & 1.88 & 4.20 & 7.20 & 7.15 & 21.5 \\
Ridge (panel) & 2.21 & 5.16 & 12.50 & 12.25 & 33.1 \\
\bottomrule
\end{tabular}
}\end{table}
\begin{table}[t]\centering\caption{Validation (2015--2019) accuracy.}\label{tab:val}
\resizebox{\columnwidth}{!}{\begin{tabular}{lrrrrr}
\toprule
Model & MASE$_1$ & MASE$_3$ & MASE$_5$ & MinT MASE$_5$ & sMAPE$_5$ \\
\midrule
GRU + emb. & 1.14 & \textbf{1.89} & \textbf{2.13} & 2.14 & 7.5 \\
TCN + emb. & 1.25 & 2.03 & 2.25 & 2.26 & 7.9 \\
LSTM + emb. & 1.08 & 1.93 & 2.54 & 2.55 & 9.3 \\
DeepAR-style & 1.27 & 2.33 & 2.76 & 2.77 & 9.9 \\
Seq2Seq-Attn & 1.62 & 2.50 & 2.81 & 2.83 & 10.1 \\
TFT-lite & 1.55 & 2.45 & 2.89 & 2.88 & 11.2 \\
ARIMA$^\dagger$ & \textbf{0.99} & 1.91 & 2.93 & 2.93 & 10.0 \\
LightGBM (panel) & 1.02 & 2.24 & 3.31 & 3.31 & 11.3 \\
RW + drift & 1.03 & 2.33 & 3.53 & 3.51 & 12.1 \\
Ridge (panel) & 1.63 & 3.73 & 5.27 & 5.20 & 18.4 \\
Naive (RW) & 1.18 & 3.29 & 5.42 & 4.99 & 20.8 \\
\bottomrule
\end{tabular}
}\end{table}
\begin{table}[t]\centering\caption{Permutation importance of feature groups ($\Delta$ val MASE).}\label{tab:imp}
\resizebox{\columnwidth}{!}{\begin{tabular}{lrrrrrrrrr}
\toprule
Group & GRU + emb. & LightGBM (panel) & lightgbm_sector & mt_gru & mt_tft & Ridge (panel) & Seq2Seq-Attn & TCN + emb. & TFT-lite \\
\midrule
climate & 0.025 & -0.002 & 0.092 & 0.054 & 0.264 & 0.197 & 0.024 & 0.066 & 0.021 \\
demography & 0.113 & 0.046 & 0.397 & 0.086 & 0.296 & 0.570 & 0.070 & 0.071 & 0.123 \\
economy & 0.328 & 0.104 & 0.921 & 0.469 & 0.024 & 1.308 & 0.177 & 0.151 & 0.105 \\
energy & 0.232 & 0.059 & nan & nan & nan & 1.658 & 0.151 & nan & -0.009 \\
history & 0.673 & 0.741 & nan & nan & nan & 0.589 & 0.483 & nan & 0.517 \\
policy & 0.130 & 0.000 & 0.000 & 0.065 & 0.005 & 0.369 & 0.104 & -0.024 & 0.006 \\
sector_drivers & 0.029 & 0.034 & 0.093 & 0.052 & 0.063 & nan & -0.049 & 0.016 & 0.053 \\
sector_history & 1.357 & 6.396 & 6.145 & 1.058 & 2.040 & nan & 1.385 & 0.967 & 1.095 \\
state_total_history & 0.020 & 0.030 & 0.459 & 0.044 & 0.144 & nan & -0.007 & 0.015 & 0.020 \\
transport & 0.296 & 0.161 & nan & nan & nan & 0.407 & -0.024 & nan & 0.595 \\
\bottomrule
\end{tabular}
}\end{table}
\todo{Discuss: COVID-19 shock dominates test h=1; attention concentrates on the most recent 4--6 input years;
history and economy groups dominate importance; decoupling (Insights page) and watchlist findings.}


\section{State$\times$Sector Forecasting (v2)}
\textbf{Data.} EDGAR 2025 CO\textsubscript{2}, CH\textsubscript{4} and N\textsubscript{2}O sub-sector grids (1,785 grids,
1990--2024) are mapped to seven sectors (Power, Transport, Industry, Buildings, Agriculture, Waste, Other) and converted to
CO\textsubscript{2}e with AR6 GWP-100 (CH\textsubscript{4} 29.8 fossil / 27.0 non-fossil; N\textsubscript{2}O 273). Other is the
residual against EDGAR CO\textsubscript{2} totals, so $\sum_k$ sector CO\textsubscript{2} equals the v1 state series exactly.
\textbf{Models.} A unit is a (state, sector) pair with embedding $e = E_{\text{state}}[s] + E_{\text{sector}}[k]$. Inputs add
lagged own-sector and state-total history and sector drivers. The multi-task model encodes each sector's window with a
shared GRU (or TFT-style variable selection + GRU + temporal attention), mixes the seven sector tokens with a
Transformer layer, and emits per-sector P10/P50/P90.
\textbf{Reconciliation.} With $S\in\{0,1\}^{296\times252}$ (national, 36 states, 7 sectors, 252 cells) and diagonal $W$ of
train residual variances, $\tilde{y}=S(S^\top W^{-1}S)^{-1}S^\top W^{-1}\hat{y}$. State totals use a direct state-total TFT
as base forecast; sector and national totals use drift.

\begin{table}[t]\centering\caption{State$\times$sector test accuracy (material cells, train mean $\ge$10 kt).}\label{tab:sector}
\resizebox{\columnwidth}{!}{\begin{tabular}{lrrrrrr}
\toprule
 & \multicolumn{3}{c}{MASE (unreconciled)} & \multicolumn{3}{c}{MASE (two-level MinT)} \\
Model & $h{=}1$ & $h{=}3$ & $h{=}5$ & $h{=}1$ & $h{=}3$ & $h{=}5$ \\
\midrule
TFT-lite & 1.44 & 2.49 & 2.28 & 1.50 & \textbf{2.61} & \textbf{2.50} \\
MT-GRU (ours) & 1.41 & 2.73 & 2.53 & 1.48 & 2.74 & 2.63 \\
TCN + emb. & 1.49 & 2.84 & 2.64 & 1.50 & 2.80 & 2.70 \\
GRU + emb. & 1.53 & 2.71 & 2.66 & 1.57 & 2.78 & 2.76 \\
MT-TFT (ours) & 1.37 & 2.72 & 3.07 & 1.43 & 2.83 & 3.20 \\
Naive (RW) & 1.69 & 3.15 & 3.65 & 1.62 & 2.90 & 3.21 \\
Seq2Seq-Attn & 1.51 & 2.88 & 3.26 & 1.55 & 2.88 & 3.26 \\
LightGBM (panel) & 1.31 & 2.58 & 3.33 & \textbf{1.38} & 2.71 & 3.42 \\
RW + drift & 1.38 & 3.00 & 3.79 & 1.44 & 3.04 & 3.77 \\
LightGBM (per sector) & 1.45 & 2.85 & 3.86 & 1.50 & 2.92 & 3.86 \\
\bottomrule
\end{tabular}
}\end{table}
\begin{table}[t]\centering\caption{Reconciled MASE at $h{=}5$ by sector.}\label{tab:bysector}
\resizebox{\columnwidth}{!}{\begin{tabular}{lrrrrrr}
\toprule
Sector & Naive (RW) & LightGBM (panel) & GRU + emb. & TFT-lite & MT-GRU (ours) & MT-TFT (ours) \\
\midrule
Agriculture & 3.29 & 2.40 & 2.28 & 1.80 & 2.14 & 1.89 \\
Buildings & 1.96 & 1.32 & 1.71 & 1.17 & 1.61 & 1.63 \\
Industry & 2.54 & 3.95 & 2.88 & 2.69 & 2.65 & 3.87 \\
Other & 3.24 & 3.15 & 3.32 & 2.83 & 2.90 & 4.12 \\
Power & 5.75 & 6.14 & 5.18 & 5.68 & 5.32 & 5.60 \\
Transport & 2.15 & 5.74 & 2.19 & 1.92 & 1.90 & 3.34 \\
Waste & 4.22 & 1.80 & 2.41 & 2.24 & 2.60 & 2.68 \\
\bottomrule
\end{tabular}
}\end{table}

\section{Experiments (v2)}
\textbf{E10 disaggregation} (Table~\ref{tab:e10}): state-total MASE of bottom-up sector sums vs a direct state-total model.
\textbf{E11 multi-task} (Table~\ref{tab:e11}). \textbf{E12 reconciliation} (Table~\ref{tab:e12}).
\textbf{E13 counterfactual sanity}: CI asserts monotonicity (more renewables never increase power emissions; EVs reduce
transport and add grid load; BAU is the identity). \textbf{E14 transfer} (Table~\ref{tab:e14}): six held-out states predicted
zero-shot with the mean state embedding. \todo{interpret once tables are final.}
\begin{table}[t]\centering\caption{E10: state totals, bottom-up vs direct.}\label{tab:e10}\resizebox{\columnwidth}{!}{\begin{tabular}{llrrr}
\toprule
Model & Variant & $h{=}1$ & $h{=}3$ & $h{=}5$ \\
\midrule
RW + drift & bottom up reconciled & 1.71 & 3.52 & 3.88 \\
RW + drift & bottom up sectors & 1.55 & 3.50 & 4.30 \\
RW + drift & direct state total & 1.99 & 3.80 & 4.24 \\
GRU + emb. & bottom up reconciled & 1.72 & 3.31 & 2.83 \\
GRU + emb. & bottom up sectors & 1.63 & 3.00 & 2.57 \\
GRU + emb. & direct state total & 1.99 & 3.80 & 4.24 \\
LightGBM (panel) & bottom up reconciled & 1.58 & 3.16 & 3.76 \\
LightGBM (panel) & bottom up sectors & 1.42 & 2.85 & 3.79 \\
LightGBM (panel) & direct state total & 1.99 & 3.80 & 4.24 \\
LightGBM (per sector) & bottom up reconciled & 1.62 & 3.36 & 4.55 \\
LightGBM (per sector) & bottom up sectors & 1.50 & 3.15 & 4.79 \\
LightGBM (per sector) & direct state total & 1.99 & 3.80 & 4.24 \\
MT-GRU (ours) & bottom up reconciled & 1.68 & 3.30 & 3.26 \\
MT-GRU (ours) & bottom up sectors & 1.47 & 3.06 & 2.87 \\
MT-GRU (ours) & direct state total & 1.99 & 3.80 & 4.24 \\
MT-TFT (ours) & bottom up reconciled & 1.77 & 3.60 & 4.11 \\
MT-TFT (ours) & bottom up sectors & 1.67 & 3.44 & 3.95 \\
MT-TFT (ours) & direct state total & 1.99 & 3.80 & 4.24 \\
Naive (RW) & bottom up reconciled & 1.83 & 3.40 & 3.10 \\
Naive (RW) & bottom up sectors & 1.98 & 3.74 & 4.27 \\
Naive (RW) & direct state total & 1.99 & 3.80 & 4.24 \\
Seq2Seq-Attn & bottom up reconciled & 1.75 & 3.37 & 3.24 \\
Seq2Seq-Attn & bottom up sectors & 1.71 & 3.31 & 3.32 \\
Seq2Seq-Attn & direct state total & 1.99 & 3.80 & 4.24 \\
TCN + emb. & bottom up reconciled & 1.74 & 3.41 & 3.44 \\
TCN + emb. & bottom up sectors & 1.68 & 3.20 & 2.78 \\
TCN + emb. & direct state total & 1.99 & 3.80 & 4.24 \\
TFT-lite & bottom up reconciled & 1.90 & 3.58 & 3.72 \\
TFT-lite & bottom up sectors & 2.03 & 3.36 & 3.18 \\
TFT-lite & direct state total & 1.99 & 3.80 & 4.24 \\
\bottomrule
\end{tabular}
}\end{table}
\begin{table}[t]\centering\caption{E11: multi-task vs per-sector heads.}\label{tab:e11}\resizebox{\columnwidth}{!}{\begin{tabular}{llrrr}
\toprule
Model & Variant & $h{=}1$ & $h{=}3$ & $h{=}5$ \\
\midrule
GRU + emb. & per sector head & 1.54 & 2.84 & 2.92 \\
MT-GRU (ours) & multi task & 1.55 & 2.90 & 2.96 \\
MT-TFT (ours) & multi task & 1.53 & 2.79 & 2.85 \\
TFT-lite & per sector head & 1.45 & 2.53 & 2.50 \\
\bottomrule
\end{tabular}
}\end{table}
\begin{table}[t]\centering\caption{E12: unreconciled vs two-level MinT.}\label{tab:e12}\resizebox{\columnwidth}{!}{\begin{tabular}{llrrr}
\toprule
Model & Variant & $h{=}1$ & $h{=}3$ & $h{=}5$ \\
\midrule
RW + drift & two level MinT & 1.58 & 3.11 & 3.87 \\
RW + drift & unreconciled & 1.42 & 3.04 & 3.95 \\
GRU + emb. & two level MinT & 1.65 & 2.83 & 2.89 \\
GRU + emb. & unreconciled & 1.54 & 2.84 & 2.92 \\
LightGBM (panel) & two level MinT & 1.48 & 2.66 & 3.40 \\
LightGBM (panel) & unreconciled & 1.36 & 2.57 & 3.45 \\
LightGBM (per sector) & two level MinT & 1.59 & 2.98 & 4.06 \\
LightGBM (per sector) & unreconciled & 1.51 & 2.95 & 4.09 \\
MT-GRU (ours) & two level MinT & 1.66 & 2.89 & 2.97 \\
MT-GRU (ours) & unreconciled & 1.55 & 2.90 & 2.96 \\
MT-TFT (ours) & two level MinT & 1.63 & 2.82 & 2.90 \\
MT-TFT (ours) & unreconciled & 1.53 & 2.79 & 2.85 \\
Naive (RW) & two level MinT & 1.78 & 2.97 & 3.19 \\
Naive (RW) & unreconciled & 1.73 & 3.17 & 3.72 \\
Seq2Seq-Attn & two level MinT & 1.80 & 2.86 & 2.98 \\
Seq2Seq-Attn & unreconciled & 1.74 & 2.81 & 2.99 \\
TCN + emb. & two level MinT & 1.60 & 2.78 & 3.10 \\
TCN + emb. & unreconciled & 1.48 & 2.74 & 3.11 \\
TFT-lite & two level MinT & 1.59 & 2.63 & 2.54 \\
TFT-lite & unreconciled & 1.45 & 2.53 & 2.50 \\
\bottomrule
\end{tabular}
}\end{table}
\begin{table}[t]\centering\caption{E14: zero-shot transfer to held-out states.}\label{tab:e14}\resizebox{\columnwidth}{!}{\begin{tabular}{llrrr}
\toprule
Model & Variant & $h{=}1$ & $h{=}3$ & $h{=}5$ \\
\midrule
GRU + emb. & held out zero shot & 1.39 & 2.35 & 2.30 \\
GRU + emb. & trained on all & 1.42 & 2.47 & 2.47 \\
MT-TFT (ours) & held out zero shot & 2.13 & 5.16 & 7.70 \\
MT-TFT (ours) & trained on all & 1.44 & 2.65 & 2.62 \\
\bottomrule
\end{tabular}
}\end{table}

\textbf{Scenario engine.} Renewable share, EV adoption and industrial policy never vary in the training window, so the
networks cannot learn their effect; we therefore apply explicit identities to the reconciled forecast (power $\times
(1-\Delta_{\text{RE}}/\text{fossil share})$; EV load at the state grid factor; compounding industry growth; building
efficiency) and widen bands by $\pm$20\% of the delta. \textbf{ATAR} generates anomaly, decoupling, watchlist,
leader/laggard, driver, comparison and trend insights from templates over the same data, each with a machine-readable
evidence record.

\section{Limitations}
EDGAR spatial allocation is proxy-based (typical sub-national uncertainty $\pm$20--50\%); CEA covers only FY2017+ and
conventional plants; GSDP before FY2011 and for Ladakh/Lakshadweep/DNH\&DD is imputed; climate uses one point per state;
city values are buffer estimates that may cross state lines. Test years include the COVID-19 shock. Sector allocation inherits EDGAR proxy uncertainty (higher for agriculture and waste); several sector drivers (MoRTH, livestock, MSW) require manual inputs; scenario levers are structural, not learned.

\bibliographystyle{IEEEtran}
\begin{thebibliography}{9}
\bibitem{edgar} M. Crippa \emph{et al.}, ``EDGAR 2025 GHG emissions,'' JRC, 2025.
\bibitem{cea} Central Electricity Authority, ``CO\textsubscript{2} Baseline Database for the Indian Power Sector, v22,'' 2026.
\bibitem{mint} S. Wickramasuriya, G. Athanasopoulos, R. Hyndman, ``Optimal forecast reconciliation for hierarchical and grouped time series through trace minimization,'' \emph{JASA}, 2019.
\bibitem{tft} B. Lim \emph{et al.}, ``Temporal Fusion Transformers for interpretable multi-horizon time series forecasting,'' \emph{IJF}, 2021.
\bibitem{deepar} D. Salinas \emph{et al.}, ``DeepAR: Probabilistic forecasting with autoregressive recurrent networks,'' \emph{IJF}, 2020.
\bibitem{oxcgrt} T. Hale \emph{et al.}, ``A global panel database of pandemic policies (OxCGRT),'' \emph{Nat. Hum. Behav.}, 2021.
\end{thebibliography}
\end{document}
